\section{Données spectrales d'une réalisation supersymétrique : \((\mathcal H_{\rm phys},H,\mathcal Q)\), supercharge et compactification en dimension \(11\)}
\label{sec:md-sucesora-terna-supercarga-m}

L'interface dimensionnelle vers la théorie M exige de distinguer la structure exacte
de départ, le relèvement opératoire et les champs physiques du codomaine. Soit
\[
 \mathfrak S_M=(\mathcal H_{\mathrm{phys}},H,Q)
\]
un triplet spectral dans lequel \(\mathcal H_{\mathrm{phys}}\) est l'espace de
Hilbert après imposition des contraintes de jauge, \(H\) est un opérateur
autoadjoint non négatif et \(Q\) est une supercharge impaire, densément définie et
compatible avec la graduation
\(\mathcal H_{\mathrm{phys}}=\mathcal H_+\oplus\mathcal H_-\). La relation
\begin{equation}
 Q:\mathcal H_+\rightleftarrows\mathcal H_-,
 \qquad
 Q^2=H
 \label{eq:md-sucesora-supercarga-cuadrado}
\end{equation}
s'interprète sur leur domaine commun. Elle ne remplace pas l'algèbre supersymétrique
relativiste ; elle est sa réduction hermitienne de dimension \(1\) et fixe la donnée minimale
que doit conserver une réalisation spectrale.

\subsection{Relèvement opératoire de l'état APP--TRIT--TPK}

Soit \(\mathcal H_{\mathrm{HMT}}\) la complétion des sections
survivantes avec phase, orientation, quotient, retenue, frontière, chemin et
mémoire, et soit
\[
 W:\mathcal H_{\mathrm{HMT}}\longrightarrow\mathcal H_{\mathrm{phys}}
\]
une isométrie de réalisation. Un relèvement opératoire admissible conserve
la composition d'APP, l'orientation de TRIT et la dynamique complète du TPK,
et entrelace
\begin{equation}
 W\widehat H=HW,
 \qquad
 W\widehat Q=QW,
 \qquad
 \widehat Q^{2}=\widehat H.
 \label{eq:md-sucesora-lift-app-supercarga}
\end{equation}
La première égalité transporte le générateur, la deuxième préserve la graduation
et la troisième empêche d'introduire \(H\) comme une donnée étrangère à l'antécédent
généalogique. Ces identités sont des obligations du pont : l'existence de
complétions fermionique et bosonique ne démontre pas, à elle seule,
\eqref{eq:md-sucesora-supercarga-cuadrado}.

\Needspace{42\baselineskip}
\subsection{Codomaine physique à onze dimensions}

La réalisation en dimension \(11\) s'exprime au moyen de
\[
 \mathfrak R_M:\mathcal X_{\mathrm{HMT}}^{\mathrm{enr}}
 \longrightarrow
 \mathcal X_{11}^{\mathrm{phys}},
\]
dont le codomaine porte la métrique \(g_{MN}\), la \(3\)-forme \(C_{MNP}\), le
gravitino \(\psi_M\), l'action et ses transformations de supersymétrie. La
compactification sur une direction doit transporter simultanément ces
champs, le produit intérieur et la graduation de \(\mathfrak S_M\). Les cordes
apparaissent sur l'écran de dimension (10) ; les branes se couplent à la
\(3\)-forme et à ses duales ; et la dualité T échange impulsion et
enroulement au moyen de
\[
 (m,w,R)\longmapsto(w,m,R^{-1})
\]
sans altérer le spectre du secteur compact.

L'action de Polyakov
\[
 S_{\mathrm P}[X,h]
 =-\frac{T}{2}\int d^2\sigma\,
 \sqrt{-h}\,h^{ab}g_{MN}(X)
 \partial_aX^M\partial_bX^N
\]
fixe la dynamique de la surface d'univers une fois déterminés \(T\),
\(h_{ab}\), les applications \(X^M\) et les conditions de frontière. La chaîne
d'écrans \(12\to11\to10\) et \(26\to24\), l'involution de rayon et la
réduction réticulaire sont des résultats algébriques antérieurs ; l'action, le
spectre et les amplitudes constituent la réalisation physique ultérieure.

\subsection{Compatibilité du triplet spectral \((\mathcal H_{\rm phys},H,\mathcal Q)\) avec le codomaine physique en dimension \(11\)}

Une réalisation est admissible si elle préserve simultanément :
\begin{enumerate}
 \item le prolongement compatible de l'état à toute profondeur ;
 \item les domaines et codomaines de \(H\), \(Q\) et des observables ;
 \item les identités d'entrelacement
       \eqref{eq:md-sucesora-lift-app-supercarga} ;
 \item l'action unitaire de la dualité T sur le spectre ;
 \item la descente des champs de dimension \(11\) vers la compactification
       déclarée ;
 \item la compatibilité avec le produit intérieur physique et les contraintes
       de jauge.
\end{enumerate}
L'échec de l'une quelconque de ces conditions distingue une analogie de rangs
d'une réalisation physique. La formulation n'identifie pas automatiquement le
corridor exceptionnel à la théorie M : les deux publications proviennent du même
état enrichi, mais utilisent des foncteurs et des codomaines distincts.
